 %Cinemática
 %Movimento retilíneo uniforme
 \newcommand{\posicaofinal}{$s = s_0 + v \times \Delta t$}
 %s: posição final (m)
 %s_0: posição inicial (m)
 %v: velocidade (m/s)
 %\Delta t: intervalo de tempo (s)
 
 %Movimento retilíneo uniformemente variado
 \newcommand{\posicaofinall}{$s = s_0 + v_0 \times t + \frac{1}{2} a \times t^2$}
 %s: posição final (m)
 %s_0: posição inicial (m)
 %v_0: velocidade inicial (m/s)
 %a: aceleração (m/s2)
 %t: tempo (s)
 
 \newcommand{\velocidadefinal}{$v = v_0 + a \times t$}
 %v: velocidade final (m/s)
 %v_0: velocidade inicial (m/s)
 %a: aceleração (m/s2)
 %t: tempo (s)
 
 \newcommand{\velocidadefinall}{$v^2 = {v_0}^2 + 2 \times a \times \Delta s$}
 %v: velocidade final (m/s)
 %v_0: velocidade inicial (m/s)
 %a: aceleração (m/s2)
 %\Delta s: distância percorrida (m)
 
 %Movimento Circular Uniforme
 \newcommand{\velocidadelinear}{$v = \omega \times R$}
 %v: velocidade (m/s)
 %\omega: velocidade angular (rad/s)
 %R: raio da curvatura da trajetória (m)
 %
 \newcommand{\periodo}{$T  = \frac{1}{f}$}
 %
 %T: período (s)
 %f: frequência (Hz)
 %
 \newcommand{\velocidadeangular}{$\omega = 2  \pi f$}
 %\omega: velocidade angular (rad/s)
 %f: frequência (Hz)
 %
 \newcommand{\aceleracaocentripeta}{$\alpha_{cp} = {V^2}{R}$}
 %
 %acp: aceleração centrípeta (m/s2)
 %v: velocidade (m/s)
 %R: raio da curvatura da trajetória (m)
 
 %%Dinâmica
 
 \newcommand{\forcaresultante}{$F_R = m \times a$}
 %
 %FR: força resultante (N)
 %m: massa (kg)
 %a: aceleração (m/s2)
 %
 \newcommand{\forcapeso}{$P = m \times g}
 %P: peso (N)
 %m: massa (kg)
 %g: aceleração da gravidade (m/s2)
 %
 \newcommand{\forcaatrito}{$f_{at} = \mu \times N$}
 %fat: força de atrito (N)
 %µ: coeficiente de atrito
 %N: força normal (N)
 %
 \newcommand{\forcaelastica}{$f_{el} = k \times x$}
 %fel: força elástica (N)
 %k: constante elástica da mola (N/m)
 %x: deformação da mola (m)
 
 %Trabalho, Energia e Potência
 
 \newcommand{\trabalho}{$T = F \times d \times \cos{\theta}$}
 %T: trabalho (J)
 %F: força (N)
 %d: deslocamento(m)
 %θ:ângulo entre a direção da força e do deslocamento
 %
 \newcommand{\energiacinetica}{$E_c = \frac{1}{2}\times m \times v^2$}
 %Ec: energia cinética (J)
 %m: massa (kg)
 %v: velocidade (m/s)
 %
 \newcommand{\energiapotencial}{$E_p = m \times g \times h$}
 %Ep: energia potencial gravitacional (J)
 %m: massa (kg)
 %g: aceleração da gravidade (m/s2)
 %h: altura (m)
 %
 \newcommand{\energiaelastica}{$E_{el} = \frac{1}{2}\times k \times x^2$}
 %Eel: energia potencial elástica (J)
 %k: constante elástica da mola (N/m)
 %x: deformação da mola (m)
 %
 \newcommand{\potencia}{$P=\frac{W}{\Delta t}$}
 
 %P: potência (w)
 %T:trabalho (J)
 %Δt: intervalo de tempo (s)
 %
 \newcommand{\quantidademovimento}{$Q = m \times v$}
 
 %Q: quantidade de movimento (kg.m/s)
 %m: massa (kg)
 %v: velocidade (m/s)
 %
 \newcommand{\impulso}{$I = F \times \Delta t$}
 
 %I: impulso (N.s)
 %F: força (N)
 %\Delta t: intervalo de tempo (s)
 %
 %Hidrostática

 %Empuxo
 %
  \newcommand{\pressao}{$p=\frac{F}{A}$}
 
 %p: pressão (N/m2)
 %F: força (N)
 %A: área (m2)
 %
  \newcommand{\densidade}{$\rho = \frac{m}{v}$}
 %
 %ρ: densidade (kg/m3)
 %m: massa (kg)
 %V: volume (m3)
 %
 %pt = patm + ρ \times g \times h
 %pt: pressão total (N/m2)
 %patm: pressão atmosférica(N/m2)
 %ρ: densidade (kg/m3)
 %g: aceleração da gravidade (m/s2)
 %h: altura (m)
 %
  \newcommand{\empuxo}{$E = \rho \times g \times V$}
 %E: empuxo (N)
 %ρ: densidade (kg/m3)
 %g: aceleração da gravidade (m/s2)
 %V: volume de líquido deslocado (m3)
 %
 %Sistema solar
 %
 \newcommand{\periodoplaneta}{$T^2 = K \times r^3$}
 %T: período do planeta (u.a)
 %K: constante de proporcionalidade
 %r: raio médio (u.a)
 %
  \newcommand{\forcagravitacional}{$F_G = G \times \frac{M_1 \times M_2}{d^2}$}
 %
 %FG: força gravitacional (N)
 %G: constante de gravitação universal (N.m2/kg2)
 %M1: massa do corpo 1 (kg)
 %M2: massa do corpo 2 (kg)
 %d: distância (m)
 %

 %Termologia
 %
 %Escalas termométricas
  %Cinemática
 %Movimento retilíneo uniforme
 \newcommand{\posicaofinal}{$s = s_0 + v \times \Delta t$}
 %s: posição final (m)
 %s_0: posição inicial (m)
 %v: velocidade (m/s)
 %\Delta t: intervalo de tempo (s)
 
 %Movimento retilíneo uniformemente variado
 \newcommand{\posicaofinall}{$s = s_0 + v_0 \times t + \frac{1}{2} a \times t^2$}
 %s: posição final (m)
 %s_0: posição inicial (m)
 %v_0: velocidade inicial (m/s)
 %a: aceleração (m/s2)
 %t: tempo (s)
 
 \newcommand{\velocidadefinal}{$v = v_0 + a \times t$}
 %v: velocidade final (m/s)
 %v_0: velocidade inicial (m/s)
 %a: aceleração (m/s2)
 %t: tempo (s)
 
 \newcommand{\velocidadefinall}{$v^2 = {v_0}^2 + 2 \times a \times \Delta s$}
 %v: velocidade final (m/s)
 %v_0: velocidade inicial (m/s)
 %a: aceleração (m/s2)
 %\Delta s: distância percorrida (m)
 
 %Movimento Circular Uniforme
 \newcommand{\velocidadelinear}{$v = \omega \times R$}
 %v: velocidade (m/s)
 %\omega: velocidade angular (rad/s)
 %R: raio da curvatura da trajetória (m)
 %
 \newcommand{\periodo}{$T  = \frac{1}{f}$}
 %
 %T: período (s)
 %f: frequência (Hz)
 %
 \newcommand{\velocidadeangular}{$\omega = 2  \pi f$}
 %\omega: velocidade angular (rad/s)
 %f: frequência (Hz)
 %
 \newcommand{\aceleracaocentripeta}{$\alpha_{cp} = {V^2}{R}$}
 %
 %acp: aceleração centrípeta (m/s2)
 %v: velocidade (m/s)
 %R: raio da curvatura da trajetória (m)
 
 %%Dinâmica
 
 \newcommand{\forcaresultante}{$F_R = m \times a$}
 %
 %FR: força resultante (N)
 %m: massa (kg)
 %a: aceleração (m/s2)
 %
 \newcommand{\forcapeso}{$P = m \times g}
 %P: peso (N)
 %m: massa (kg)
 %g: aceleração da gravidade (m/s2)
 %
 \newcommand{\forcaatrito}{$f_{at} = \mu \times N$}
 %fat: força de atrito (N)
 %µ: coeficiente de atrito
 %N: força normal (N)
 %
 \newcommand{\forcaelastica}{$f_{el} = k \times x$}
 %fel: força elástica (N)
 %k: constante elástica da mola (N/m)
 %x: deformação da mola (m)
 
 %Trabalho, Energia e Potência
 
 \newcommand{\trabalho}{$T = F \times d \times \cos{\theta}$}
 %T: trabalho (J)
 %F: força (N)
 %d: deslocamento(m)
 %θ:ângulo entre a direção da força e do deslocamento
 %
 \newcommand{\energiacinetica}{$E_c = \frac{1}{2}\times m \times v^2$}
 %Ec: energia cinética (J)
 %m: massa (kg)
 %v: velocidade (m/s)
 %
 \newcommand{\energiapotencial}{$E_p = m \times g \times h$}
 %Ep: energia potencial gravitacional (J)
 %m: massa (kg)
 %g: aceleração da gravidade (m/s2)
 %h: altura (m)
 %
 \newcommand{\energiaelastica}{$E_{el} = \frac{1}{2}\times k \times x^2$}
 %Eel: energia potencial elástica (J)
 %k: constante elástica da mola (N/m)
 %x: deformação da mola (m)
 %
 \newcommand{\potencia}{$P=\frac{W}{\Delta t}$}
 
 %P: potência (w)
 %T:trabalho (J)
 %Δt: intervalo de tempo (s)
 %
 \newcommand{\quantidademovimento}{$Q = m \times v$}
 
 %Q: quantidade de movimento (kg.m/s)
 %m: massa (kg)
 %v: velocidade (m/s)
 %
 \newcommand{\impulso}{$I = F \times \Delta t$}
 
 %I: impulso (N.s)
 %F: força (N)
 %\Delta t: intervalo de tempo (s)
 %
 %Hidrostática

 %Empuxo
 %
  \newcommand{\pressao}{$p=\frac{F}{A}$}
 
 %p: pressão (N/m2)
 %F: força (N)
 %A: área (m2)
 %
  \newcommand{\densidade}{$\rho = \frac{m}{v}$}
 %
 %ρ: densidade (kg/m3)
 %m: massa (kg)
 %V: volume (m3)
 %
 %pt = patm + ρ \times g \times h
 %pt: pressão total (N/m2)
 %patm: pressão atmosférica(N/m2)
 %ρ: densidade (kg/m3)
 %g: aceleração da gravidade (m/s2)
 %h: altura (m)
 %
  \newcommand{\empuxo}{$E = \rho \times g \times V$}
 %E: empuxo (N)
 %ρ: densidade (kg/m3)
 %g: aceleração da gravidade (m/s2)
 %V: volume de líquido deslocado (m3)
 %
 %Sistema solar
 %
 \newcommand{\periodoplaneta}{$T^2 = K \times r^3$}
 %T: período do planeta (u.a)
 %K: constante de proporcionalidade
 %r: raio médio (u.a)
 %
  \newcommand{\forcagravitacional}{$F_G = G \times \frac{M_1 \times M_2}{d^2}$}
 %
 %FG: força gravitacional (N)
 %G: constante de gravitação universal (N.m2/kg2)
 %M1: massa do corpo 1 (kg)
 %M2: massa do corpo 2 (kg)
 %d: distância (m)
 %

 %Termologia
 %
 %Escalas termométricas
  %Cinemática
 %Movimento retilíneo uniforme
 \newcommand{\posicaofinal}{$s = s_0 + v \times \Delta t$}
 %s: posição final (m)
 %s_0: posição inicial (m)
 %v: velocidade (m/s)
 %\Delta t: intervalo de tempo (s)
 
 %Movimento retilíneo uniformemente variado
 \newcommand{\posicaofinall}{$s = s_0 + v_0 \times t + \frac{1}{2} a \times t^2$}
 %s: posição final (m)
 %s_0: posição inicial (m)
 %v_0: velocidade inicial (m/s)
 %a: aceleração (m/s2)
 %t: tempo (s)
 
 \newcommand{\velocidadefinal}{$v = v_0 + a \times t$}
 %v: velocidade final (m/s)
 %v_0: velocidade inicial (m/s)
 %a: aceleração (m/s2)
 %t: tempo (s)
 
 \newcommand{\velocidadefinall}{$v^2 = {v_0}^2 + 2 \times a \times \Delta s$}
 %v: velocidade final (m/s)
 %v_0: velocidade inicial (m/s)
 %a: aceleração (m/s2)
 %\Delta s: distância percorrida (m)
 
 %Movimento Circular Uniforme
 \newcommand{\velocidadelinear}{$v = \omega \times R$}
 %v: velocidade (m/s)
 %\omega: velocidade angular (rad/s)
 %R: raio da curvatura da trajetória (m)
 %
 \newcommand{\periodo}{$T  = \frac{1}{f}$}
 %
 %T: período (s)
 %f: frequência (Hz)
 %
 \newcommand{\velocidadeangular}{$\omega = 2  \pi f$}
 %\omega: velocidade angular (rad/s)
 %f: frequência (Hz)
 %
 \newcommand{\aceleracaocentripeta}{$\alpha_{cp} = {V^2}{R}$}
 %
 %acp: aceleração centrípeta (m/s2)
 %v: velocidade (m/s)
 %R: raio da curvatura da trajetória (m)
 
 %%Dinâmica
 
 \newcommand{\forcaresultante}{$F_R = m \times a$}
 %
 %FR: força resultante (N)
 %m: massa (kg)
 %a: aceleração (m/s2)
 %
 \newcommand{\forcapeso}{$P = m \times g}
 %P: peso (N)
 %m: massa (kg)
 %g: aceleração da gravidade (m/s2)
 %
 \newcommand{\forcaatrito}{$f_{at} = \mu \times N$}
 %fat: força de atrito (N)
 %µ: coeficiente de atrito
 %N: força normal (N)
 %
 \newcommand{\forcaelastica}{$f_{el} = k \times x$}
 %fel: força elástica (N)
 %k: constante elástica da mola (N/m)
 %x: deformação da mola (m)
 
 %Trabalho, Energia e Potência
 
 \newcommand{\trabalho}{$T = F \times d \times \cos{\theta}$}
 %T: trabalho (J)
 %F: força (N)
 %d: deslocamento(m)
 %θ:ângulo entre a direção da força e do deslocamento
 %
 \newcommand{\energiacinetica}{$E_c = \frac{1}{2}\times m \times v^2$}
 %Ec: energia cinética (J)
 %m: massa (kg)
 %v: velocidade (m/s)
 %
 \newcommand{\energiapotencial}{$E_p = m \times g \times h$}
 %Ep: energia potencial gravitacional (J)
 %m: massa (kg)
 %g: aceleração da gravidade (m/s2)
 %h: altura (m)
 %
 \newcommand{\energiaelastica}{$E_{el} = \frac{1}{2}\times k \times x^2$}
 %Eel: energia potencial elástica (J)
 %k: constante elástica da mola (N/m)
 %x: deformação da mola (m)
 %
 \newcommand{\potencia}{$P=\frac{W}{\Delta t}$}
 
 %P: potência (w)
 %T:trabalho (J)
 %Δt: intervalo de tempo (s)
 %
 \newcommand{\quantidademovimento}{$Q = m \times v$}
 
 %Q: quantidade de movimento (kg.m/s)
 %m: massa (kg)
 %v: velocidade (m/s)
 %
 \newcommand{\impulso}{$I = F \times \Delta t$}
 
 %I: impulso (N.s)
 %F: força (N)
 %\Delta t: intervalo de tempo (s)
 %
 %Hidrostática

 %Empuxo
 %
  \newcommand{\pressao}{$p=\frac{F}{A}$}
 
 %p: pressão (N/m2)
 %F: força (N)
 %A: área (m2)
 %
  \newcommand{\densidade}{$\rho = \frac{m}{v}$}
 %
 %ρ: densidade (kg/m3)
 %m: massa (kg)
 %V: volume (m3)
 %
 %pt = patm + ρ \times g \times h
 %pt: pressão total (N/m2)
 %patm: pressão atmosférica(N/m2)
 %ρ: densidade (kg/m3)
 %g: aceleração da gravidade (m/s2)
 %h: altura (m)
 %
  \newcommand{\empuxo}{$E = \rho \times g \times V$}
 %E: empuxo (N)
 %ρ: densidade (kg/m3)
 %g: aceleração da gravidade (m/s2)
 %V: volume de líquido deslocado (m3)
 %
 %Sistema solar
 %
 \newcommand{\periodoplaneta}{$T^2 = K \times r^3$}
 %T: período do planeta (u.a)
 %K: constante de proporcionalidade
 %r: raio médio (u.a)
 %
 \newcommand{\forcagravitacional}{$F_G = G \times \frac{M_1 \times M_2}{d^2}$}
 %
 %FG: força gravitacional (N)
 %G: constante de gravitação universal (N.m2/kg2)
 %M1: massa do corpo 1 (kg)
 %M2: massa do corpo 2 (kg)
 %d: distância (m)
 %

 %Termologia
 %
 %Escalas termométricas
 \newcommand{\escalascelcius}{$\frac{T_C}{5}=\frac{T_F-32}{9}$}
 
 %TC: temperatura em graus Celsius (ºC)
 %TF: temperatura em Fahrenheit (ºF)
 %
 \newcommand{\escalakelvin}{$T_K = T_C + 273$}
 %TK: temperatura em Kelvin (K)
 %TC: temperatura em Celsius (ºC)
 %
 %Dilatação Térmica
 \newcommand{\dilatacaolinear}{$\Delta L = L_0 \times \alpha \times \Delta t$}
 %∆L: dilatação linear (m)
 %L0: comprimento inicial (m)
 %α: coeficiente de dilatação linear (ºC-1)
 %\Delta t: variação de temperatura (ºC)
 %
 \newcommand{\dilatacaosuperficial}{$\Delta A = A_0 \times \beta \times \Delta t$}
 %∆A: dilatação superficial (m2)
 %A0: área inicial
 %β: coeficiente de dilatação superficial (ºC-1)
 %\Delta t: variação de temperatura (ºC)
 %
 \newcommand{\dilatacaovolumetrica}{$\Delta V = V_0 \times \gamma \times \Delta t$}
 %∆V: dilatação volumétrica (m3)
 %v_0: volume inicial (m3)
 %ϒ: coeficiente de dilatação volumétrico (ºC-1)
 %\Delta t: variação de temperatura (ºC)
 %
 %Calorimetria
 \newcommand{\capacidadetermica}{$C = m \times c$}
 %C: capacidade térmica (cal/ºC)*
 %m: massa (g)
 %c: calor específico (cal/gºC)*
 %
 \newcommand{\calorsensivel}{$Q = m \times c \times \Delta t$}
 %Q: quantidade de calor sensível (cal)*
 %m: massa (g)
 %c: calor específico (cal/g ºC)*
 %\Delta t: variação de temperatura (ºC)
 %
 \newcommand{\calorlatente}{$Q = m \times L$}
 %Q: quantidade de calor latente(cal)*
 %m: massa (g)
 %L: calor latente - mudança de fase (cal/g)*
 %
 %* Essas unidades não são do Sistema Internacional de Unidades
 %
 %Termodinâmica
 \newcommand{\energiainterna}{$\Delta U = Q - T$}
 %∆U: variação de energia interna (J)
 %Q: quantidade de calor (J)
 %T: trabalho (J)
 %
 $T = Q_q - Q_f$
 %T: trabalho (J)
 %Qq: quantidade de calor absorvida da fonte quente (J)
 %Qf: quantidade de calor cedida a fonte fria (J)
 %
 $R = \frac{T}{Q_q}$
 %
 %R: rendimento de uma máquina térmica
 %T: trabalho (J)
 %Qq: quantidade de calor absorvida da fonte quente (J)
 %
 $\Delta S = \frac{\Delta Q}{T}$
 %
 %\Delta s: variação de entropia (J/K)
 %∆Q: Quantidade de calor (J)
 %T: temperatura absoluta (K)
 %
 %ondas
 %
 %Velocidade de Propagação das Ondas
 $v = \lambda \times f$
 %v: velocidade de propagação de uma onda (m/s)
 %ƛ: comprimento de onda (m)
 %f: frequência (Hz)
 %
 %Espelhos Esféricos
 $\frac{1}{f}=\frac{1}{p}+\frac{1}{p'}$
 %
 %f: distância focal (cm ou m)
 %p: distância do vértice do espelho ao objeto (cm ou m)
 %p': distância do vértice do espelho a imagem (cm ou m)
 %
 $A=\frac{i}{o}+\frac{p'}{p}$
 %
 %A: aumento linear transversal
 %i: tamanho da imagem (cm ou m)
 %o: tamanho do objeto (cm ou m)
 %p': distância do vértice do espelho a imagem (cm ou m)
 %p: distância do vértice do espelho ao objeto (cm ou m)
 %
 %Refração
 $n_1 \times \sen{\theta_1} = n_2 \times \sen{\theta_2}$
 %n1: índice de refração do meio 1
 %θ1: ângulo de incidência
 %n2: índice de refração do meio 2
 %θ2: ângulo de refração
 %
 %
 %Eletrostática
 %
 $F_e=k\times \frac{q_1\times q_2}{d^2}$
 %
 %Fe: força eletrostática (N)
 %k: constante eletrostática (N.m2/C2)
 %Q1: módulo da carga 1 (C)
 %Q2: módulo da carga 2 (C)
 %d: distância entre as cargas (m)
 %
 $F = q \times E$
 %F: força eletrostática (N)
 %q: carga de prova (C)
 %E: campo elétrico (N/C)
 %
 $V=k\times \frac{Q}{d}$
 %
 %V: potencial elétrico (V)
 %k: constante eletrostática (N.m2/C2)
 %Q: carga elétrica (C)
 %d: distância (m)
 %
 %Eletricidade
 $U = R \times i$
 %U: diferença de potencial (V)
 %R: resistência elétrica (\omega)
 %i: corrente (A)
 %
 $P = U \times i$
 %P: potência elétrica (W)
 %U: diferença de potencial (V)
 %i: corrente (A)
 %
 $P = R \times i^2$
 %
 %P: potência efeito Joule (J)
 %R: resistência elétrica (\omega)
 %i: corrente (A)
 %
 $E = P \times \Delta t$
 %E: energia elétrica (J ou kWh)
 %P: potência (J ou kW)
 %\Delta t: intervalo de tempo (s ou h)
 %
 %Associação de Resistores em Série
 %Re = R1 + R2 + ...+ Rn
 %Re: resistência equivalente (\omega)
 %R1: resistência 1 (\omega)
 %R2: resistência 2 (\omega)
 %Rn: resistência n (\omega)
 %
 %Associação de Resistores em Paralelo
 % 1 sobre  R com  e subscrito  igual a  1 sobre  R com  1 subscrito  mais  1 sobre  R com  2 subscrito  mais . . .  mais  1 sobre  R com  n subscrito
 %
 %Re: resistência equivalente (\omega)
 %R1: resistência 1 (\omega)
 %R2: resistência 2 (\omega)
 %Rn: resistência n (\omega)
 %
 %Capacitores
 $C =\frac{Q}{U}$
 %
 %C: capacitância (F)
 %Q: carga elétrica (C)
 %U: diferença de potencial (V)
 %
 %
 $F_m= B \times  q  \times v \times \sen{\theta}$
 %Fm: força magnética (N)
 %B: vetor indução magnética (T)
 %| q |: módulo da carga (C)
 %v: velocidade (m/s)
 %θ: ângulo entre vetor B e a velocidade
 %
 $F_m= B \times i \times l \times \sen{theta}$
 %Fm: força magnética (N)
 %B: vetor indução magnética (T)
 %i: corrente (A)
 %l: comprimento do fio (m/s)
 %θ: ângulo entre vetor B e a corrente
 %
 $\phi = B \times A \times \cos{\theta}$
 %φ: fluxo magnético (Wb)
 %B: vetor indução magnética (T)
 %A: Área (m2)
 %θ: ângulo entre vetor B e o vetor normal a superfície da espira
 %
 $\varepsilon = \frac{\Delta \Phi}{\Delta t}$
 %
 %ε: fem induzida (V)
 %∆φ: variação do fluxo magnético (Wb)
 %\Delta t: intervalo de tempo (s)$\frac{T_C}{5}=\frac{T_F-32}{9}$
 
 %TC: temperatura em graus Celsius (ºC)
 %TF: temperatura em Fahrenheit (ºF)
 %
 $T_K = T_C + 273$
 %TK: temperatura em Kelvin (K)
 %TC: temperatura em Celsius (ºC)
 %
 %Dilatação Térmica
$\Delta L = L_0 \times \alpha \times \Delta t$
 %∆L: dilatação linear (m)
 %L0: comprimento inicial (m)
 %α: coeficiente de dilatação linear (ºC-1)
 %\Delta t: variação de temperatura (ºC)
 %
 $\Delta A = A_0 \times \beta \times \Delta t$
 %∆A: dilatação superficial (m2)
 %A0: área inicial
 %β: coeficiente de dilatação superficial (ºC-1)
 %\Delta t: variação de temperatura (ºC)
 %
 $\Delta V = V_0 \times \gamma \times \Delta t$
 %∆V: dilatação volumétrica (m3)
 %v_0: volume inicial (m3)
 %ϒ: coeficiente de dilatação volumétrico (ºC-1)
 %\Delta t: variação de temperatura (ºC)
 %
 %Calorimetria
 $C = m \times c$
 %C: capacidade térmica (cal/ºC)*
 %m: massa (g)
 %c: calor específico (cal/gºC)*
 %
 $Q = m \times c \times \Delta t$
 %Q: quantidade de calor sensível (cal)*
 %m: massa (g)
 %c: calor específico (cal/g ºC)*
 %\Delta t: variação de temperatura (ºC)
 %
 $Q = m \times L$
 %Q: quantidade de calor latente(cal)*
 %m: massa (g)
 %L: calor latente - mudança de fase (cal/g)*
 %
 %* Essas unidades não são do Sistema Internacional de Unidades
 %
 %Termodinâmica
 $\Delta U = Q - T$
 %∆U: variação de energia interna (J)
 %Q: quantidade de calor (J)
 %T: trabalho (J)
 %
 $T = Q_q - Q_f$
 %T: trabalho (J)
 %Qq: quantidade de calor absorvida da fonte quente (J)
 %Qf: quantidade de calor cedida a fonte fria (J)
 %
 $R = \frac{T}{Q_q}$
 %
 %R: rendimento de uma máquina térmica
 %T: trabalho (J)
 %Qq: quantidade de calor absorvida da fonte quente (J)
 %
 $\Delta S = \frac{\Delta Q}{T}$
 %
 %\Delta s: variação de entropia (J/K)
 %∆Q: Quantidade de calor (J)
 %T: temperatura absoluta (K)
 %
 %ondas
 %
 %Velocidade de Propagação das Ondas
 $v = \lambda \times f$
 %v: velocidade de propagação de uma onda (m/s)
 %ƛ: comprimento de onda (m)
 %f: frequência (Hz)
 %
 %Espelhos Esféricos
 $\frac{1}{f}=\frac{1}{p}+\frac{1}{p'}$
 %
 %f: distância focal (cm ou m)
 %p: distância do vértice do espelho ao objeto (cm ou m)
 %p': distância do vértice do espelho a imagem (cm ou m)
 %
 $A=\frac{i}{o}+\frac{p'}{p}$
 %
 %A: aumento linear transversal
 %i: tamanho da imagem (cm ou m)
 %o: tamanho do objeto (cm ou m)
 %p': distância do vértice do espelho a imagem (cm ou m)
 %p: distância do vértice do espelho ao objeto (cm ou m)
 %
 %Refração
 $n_1 \times \sen{\theta_1} = n_2 \times \sen{\theta_2}$
 %n1: índice de refração do meio 1
 %θ1: ângulo de incidência
 %n2: índice de refração do meio 2
 %θ2: ângulo de refração
 %
 %
 %Eletrostática
 %
 $F_e=k\times \frac{q_1\times q_2}{d^2}$
 %
 %Fe: força eletrostática (N)
 %k: constante eletrostática (N.m2/C2)
 %Q1: módulo da carga 1 (C)
 %Q2: módulo da carga 2 (C)
 %d: distância entre as cargas (m)
 %
 $F = q \times E$
 %F: força eletrostática (N)
 %q: carga de prova (C)
 %E: campo elétrico (N/C)
 %
 $V=k\times \frac{Q}{d}$
 %
 %V: potencial elétrico (V)
 %k: constante eletrostática (N.m2/C2)
 %Q: carga elétrica (C)
 %d: distância (m)
 %
 %Eletricidade
 $U = R \times i$
 %U: diferença de potencial (V)
 %R: resistência elétrica (\omega)
 %i: corrente (A)
 %
 $P = U \times i$
 %P: potência elétrica (W)
 %U: diferença de potencial (V)
 %i: corrente (A)
 %
 $P = R \times i^2$
 %
 %P: potência efeito Joule (J)
 %R: resistência elétrica (\omega)
 %i: corrente (A)
 %
 $E = P \times \Delta t$
 %E: energia elétrica (J ou kWh)
 %P: potência (J ou kW)
 %\Delta t: intervalo de tempo (s ou h)
 %
 %Associação de Resistores em Série
 %Re = R1 + R2 + ...+ Rn
 %Re: resistência equivalente (\omega)
 %R1: resistência 1 (\omega)
 %R2: resistência 2 (\omega)
 %Rn: resistência n (\omega)
 %
 %Associação de Resistores em Paralelo
 % 1 sobre  R com  e subscrito  igual a  1 sobre  R com  1 subscrito  mais  1 sobre  R com  2 subscrito  mais . . .  mais  1 sobre  R com  n subscrito
 %
 %Re: resistência equivalente (\omega)
 %R1: resistência 1 (\omega)
 %R2: resistência 2 (\omega)
 %Rn: resistência n (\omega)
 %
 %Capacitores
 $C =\frac{Q}{U}$
 %
 %C: capacitância (F)
 %Q: carga elétrica (C)
 %U: diferença de potencial (V)
 %
 %
 $F_m= B \times  q  \times v \times \sen{\theta}$
 %Fm: força magnética (N)
 %B: vetor indução magnética (T)
 %| q |: módulo da carga (C)
 %v: velocidade (m/s)
 %θ: ângulo entre vetor B e a velocidade
 %
 $F_m= B \times i \times l \times \sen{theta}$
 %Fm: força magnética (N)
 %B: vetor indução magnética (T)
 %i: corrente (A)
 %l: comprimento do fio (m/s)
 %θ: ângulo entre vetor B e a corrente
 %
 $\phi = B \times A \times \cos{\theta}$
 %φ: fluxo magnético (Wb)
 %B: vetor indução magnética (T)
 %A: Área (m2)
 %θ: ângulo entre vetor B e o vetor normal a superfície da espira
 %
 $\varepsilon = \frac{\Delta \Phi}{\Delta t}$
 %
 %ε: fem induzida (V)
 %∆φ: variação do fluxo magnético (Wb)
 %\Delta t: intervalo de tempo (s)$\frac{T_C}{5}=\frac{T_F-32}{9}$
 
 %TC: temperatura em graus Celsius (ºC)
 %TF: temperatura em Fahrenheit (ºF)
 %
 $T_K = T_C + 273$
 %TK: temperatura em Kelvin (K)
 %TC: temperatura em Celsius (ºC)
 %
 %Dilatação Térmica
$\Delta L = L_0 \times \alpha \times \Delta t$
 %∆L: dilatação linear (m)
 %L0: comprimento inicial (m)
 %α: coeficiente de dilatação linear (ºC-1)
 %\Delta t: variação de temperatura (ºC)
 %
 $\Delta A = A_0 \times \beta \times \Delta t$
 %∆A: dilatação superficial (m2)
 %A0: área inicial
 %β: coeficiente de dilatação superficial (ºC-1)
 %\Delta t: variação de temperatura (ºC)
 %
 $\Delta V = V_0 \times \gamma \times \Delta t$
 %∆V: dilatação volumétrica (m3)
 %v_0: volume inicial (m3)
 %ϒ: coeficiente de dilatação volumétrico (ºC-1)
 %\Delta t: variação de temperatura (ºC)
 %
 %Calorimetria
 $C = m \times c$
 %C: capacidade térmica (cal/ºC)*
 %m: massa (g)
 %c: calor específico (cal/gºC)*
 %
 $Q = m \times c \times \Delta t$
 %Q: quantidade de calor sensível (cal)*
 %m: massa (g)
 %c: calor específico (cal/g ºC)*
 %\Delta t: variação de temperatura (ºC)
 %
 $Q = m \times L$
 %Q: quantidade de calor latente(cal)*
 %m: massa (g)
 %L: calor latente - mudança de fase (cal/g)*
 %
 %* Essas unidades não são do Sistema Internacional de Unidades
 %
 %Termodinâmica
 $\Delta U = Q - T$
 %∆U: variação de energia interna (J)
 %Q: quantidade de calor (J)
 %T: trabalho (J)
 %
 $T = Q_q - Q_f$
 %T: trabalho (J)
 %Qq: quantidade de calor absorvida da fonte quente (J)
 %Qf: quantidade de calor cedida a fonte fria (J)
 %
 $R = \frac{T}{Q_q}$
 %
 %R: rendimento de uma máquina térmica
 %T: trabalho (J)
 %Qq: quantidade de calor absorvida da fonte quente (J)
 %
 $\Delta S = \frac{\Delta Q}{T}$
 %
 %\Delta s: variação de entropia (J/K)
 %∆Q: Quantidade de calor (J)
 %T: temperatura absoluta (K)
 %
 %ondas
 %
 %Velocidade de Propagação das Ondas
 $v = \lambda \times f$
 %v: velocidade de propagação de uma onda (m/s)
 %ƛ: comprimento de onda (m)
 %f: frequência (Hz)
 %
 %Espelhos Esféricos
 $\frac{1}{f}=\frac{1}{p}+\frac{1}{p'}$
 %
 %f: distância focal (cm ou m)
 %p: distância do vértice do espelho ao objeto (cm ou m)
 %p': distância do vértice do espelho a imagem (cm ou m)
 %
 $A=\frac{i}{o}+\frac{p'}{p}$
 %
 %A: aumento linear transversal
 %i: tamanho da imagem (cm ou m)
 %o: tamanho do objeto (cm ou m)
 %p': distância do vértice do espelho a imagem (cm ou m)
 %p: distância do vértice do espelho ao objeto (cm ou m)
 %
 %Refração
 $n_1 \times \sen{\theta_1} = n_2 \times \sen{\theta_2}$
 %n1: índice de refração do meio 1
 %θ1: ângulo de incidência
 %n2: índice de refração do meio 2
 %θ2: ângulo de refração
 %
 %
 %Eletrostática
 %
 $F_e=k\times \frac{q_1\times q_2}{d^2}$
 %
 %Fe: força eletrostática (N)
 %k: constante eletrostática (N.m2/C2)
 %Q1: módulo da carga 1 (C)
 %Q2: módulo da carga 2 (C)
 %d: distância entre as cargas (m)
 %
 $F = q \times E$
 %F: força eletrostática (N)
 %q: carga de prova (C)
 %E: campo elétrico (N/C)
 %
 $V=k\times \frac{Q}{d}$
 %
 %V: potencial elétrico (V)
 %k: constante eletrostática (N.m2/C2)
 %Q: carga elétrica (C)
 %d: distância (m)
 %
 %Eletricidade
 $U = R \times i$
 %U: diferença de potencial (V)
 %R: resistência elétrica (\omega)
 %i: corrente (A)
 %
 $P = U \times i$
 %P: potência elétrica (W)
 %U: diferença de potencial (V)
 %i: corrente (A)
 %
 $P = R \times i^2$
 %
 %P: potência efeito Joule (J)
 %R: resistência elétrica (\omega)
 %i: corrente (A)
 %
 $E = P \times \Delta t$
 %E: energia elétrica (J ou kWh)
 %P: potência (J ou kW)
 %\Delta t: intervalo de tempo (s ou h)
 %
 %Associação de Resistores em Série
 %Re = R1 + R2 + ...+ Rn
 %Re: resistência equivalente (\omega)
 %R1: resistência 1 (\omega)
 %R2: resistência 2 (\omega)
 %Rn: resistência n (\omega)
 %
 %Associação de Resistores em Paralelo
 % 1 sobre  R com  e subscrito  igual a  1 sobre  R com  1 subscrito  mais  1 sobre  R com  2 subscrito  mais . . .  mais  1 sobre  R com  n subscrito
 %
 %Re: resistência equivalente (\omega)
 %R1: resistência 1 (\omega)
 %R2: resistência 2 (\omega)
 %Rn: resistência n (\omega)
 %
 %Capacitores
 $C =\frac{Q}{U}$
 %
 %C: capacitância (F)
 %Q: carga elétrica (C)
 %U: diferença de potencial (V)
 %
 %
 $F_m= B \times  q  \times v \times \sen{\theta}$
 %Fm: força magnética (N)
 %B: vetor indução magnética (T)
 %| q |: módulo da carga (C)
 %v: velocidade (m/s)
 %θ: ângulo entre vetor B e a velocidade
 %
 $F_m= B \times i \times l \times \sen{theta}$
 %Fm: força magnética (N)
 %B: vetor indução magnética (T)
 %i: corrente (A)
 %l: comprimento do fio (m/s)
 %θ: ângulo entre vetor B e a corrente
 %
 $\phi = B \times A \times \cos{\theta}$
 %φ: fluxo magnético (Wb)
 %B: vetor indução magnética (T)
 %A: Área (m2)
 %θ: ângulo entre vetor B e o vetor normal a superfície da espira
 %
 $\varepsilon = \frac{\Delta \Phi}{\Delta t}$
 %
 %ε: fem induzida (V)
 %∆φ: variação do fluxo magnético (Wb)
 %\Delta t: intervalo de tempo (s)